\renewcommand{\footerfilename}{PostInstall.tex}

\section{Post Install, Use the Tools!}

But first, how do you get out something you started.

To exit PyRAF use the \dhl{.ex} as the command. To exit ecl use \dhl{logout} as the command.

When in a PyRAF Graphics window (popped out green window with black background), usually
a Capital \llbox{I} will escape. 

As a next-to-the last resort, use \mousebox{Ctrl-Z} to pause the window; in the bash shell
use \llbox{j} to list the background jobs (they are numbered); then enter \dhl{kill \%<n>}
and this will kill the job.

Sometimes these tasks hang so hard, that it requires finding the PID. Use the Unix command:
\begingroup \fontsize{10pt}{10pt}
\selectfont
\begin{verbatim} 
ps aux | grep -i pyraf
\end{verbatim}
\endgroup

Locate the PID and use a \dhl{kill -9 <PID>} to murder the task outright.

IRAF tasks cl and ecl are designed to run in a \dhl{xterm} or
\dhl{xgterm}, PyRAF can run from a normal terminal.  The ``cl'' or
``ecl'' commands are used to interact with IRAF in pure terms. In
these modern times, it is possible to run \dhl{cl} from within PyRAF
in a rather fluid way.

There is no real pyraf script. 

Including cl and ecl is intended to support legacy code and to support
tedious/specialized scripting from within PyRAF.

PyFAF ``is'' Python with modified command-line reader (readline). It
pre-loads a python module ``iraf.py''. The command-line command
``imstat'' has a corresponding wrapped IRAF command
iraf.imstat(...). This allows you to use IRAF within python loops.  It
is a bit messy, but crib-sheets that you develop assist. Yes it is
messy to begin with and hard to live without later on!

{\color{darkred}\begingroup \fontsize{10pt}{10pt}
\selectfont
\begin{verbatim} 
export MYVENVS="$HOME/MyEnves/"
function mypyraf { 
   source $MYVENVS/mypyrafvenv/bin/activate;  # start the venv tune to fit.
   export PATH="$PATH:/usr/local/bin";             # add /usr/local/bin
   if ! [[ "$PATH" =~ "$HOME/bin" ]] ;  then       # add my local bin
       export PATH="$PATH:$HOME/bin";
   fi
   if ! [[ "$PATH" =~ "local/astrometry/bin" ; then # add astrometry if I want to
       export PATH="$PATH:/usr/local/astrometry/bin";
   fi
   export PATH="$PATH:/usr/local/astrometry/bin";  # include astrometry.net as needed
   # Add my pyraflogin3 package to search path; make sure we get the env path too.XS
   export PYTHONPATH="$HOME/.iraf:$HOME/MyEnves/mypyrafvenv/lib/python3.10/site-packages/pyraf";
   pyraf -s;
   }
\end{verbatim}
\endgroup
}
